\documentclass[10pt,a4paper]{amsart}
\usepackage{amsmath,amssymb,amsthm,graphicx}
% For commutative diagrams we use Xy-pic
\usepackage[all]{xy}

% We use 2cell for 2-commutative diagrams.
\xyoption{2cell}
% Theorem environments.
%
\theoremstyle{plain}
\newtheorem{theorem}[subsection]{Theorem}
\newtheorem{proposition}[subsection]{Proposition}
\newtheorem{lemma}[subsection]{Lemma}

\theoremstyle{definition}
\newtheorem{definition}[subsection]{Definition}
\newtheorem{example}[subsection]{Example}
\newtheorem{exercise}[subsection]{Exercise}
\newtheorem{situation}[subsection]{Situation}

\theoremstyle{remark}
\newtheorem{remark}[subsection]{Remark}
\newtheorem{remarks}[subsection]{Remarks}

\numberwithin{equation}{subsection}

% Macros
%
\def\lim{\mathop{\mathrm{lim}}\nolimits}
\def\colim{\mathop{\mathrm{colim}}\nolimits}
\def\Spec{\mathop{\mathrm{Spec}}}
\def\Hom{\mathop{\mathrm{Hom}}\nolimits}
\def\Ext{\mathop{\mathrm{Ext}}\nolimits}
\def\SheafHom{\mathop{\mathcal{H}\!\mathit{om}}\nolimits}
\def\SheafExt{\mathop{\mathcal{E}\!\mathit{xt}}\nolimits}
\def\Sch{\mathit{Sch}}
\def\Mor{\mathop{\mathrm{Mor}}\nolimits}
\def\Ob{\mathop{\mathrm{Ob}}\nolimits}
\def\Sh{\mathop{\mathit{Sh}}\nolimits}
\def\NL{\mathop{N\!L}\nolimits}
\def\CH{\mathop{\mathrm{CH}}\nolimits}
\def\proetale{{pro\text{-}\acute{e}tale}}
\def\etale{{\acute{e}tale}}
\def\QCoh{\mathit{QCoh}}
\def\Ker{\mathop{\mathrm{Ker}}}
\def\Im{\mathop{\mathrm{Im}}}
\def\Coker{\mathop{\mathrm{Coker}}}
\def\Coim{\mathop{\mathrm{Coim}}}

% Boxtimes
%
\DeclareMathSymbol{\boxtimes}{\mathbin}{AMSa}{"02}

%
% Macros for moduli stacks/spaces
%
\def\QCohstack{\mathcal{QC}\!\mathit{oh}}
\def\Cohstack{\mathcal{C}\!\mathit{oh}}
\def\Spacesstack{\mathcal{S}\!\mathit{paces}}
\def\Quotfunctor{\mathrm{Quot}}
\def\Hilbfunctor{\mathrm{Hilb}}
\def\Curvesstack{\mathcal{C}\!\mathit{urves}}
\def\Polarizedstack{\mathcal{P}\!\mathit{olarized}}
\def\Complexesstack{\mathcal{C}\!\mathit{omplexes}}
% \Pic is the operator that assigns to X its picard group, usage \Pic(X)
% \Picardstack_{X/B} denotes the Picard stack of X over B
% \Picardfunctor_{X/B} denotes the Picard functor of X over B
\def\Pic{\mathop{\mathrm{Pic}}\nolimits}
\def\Picardstack{\mathcal{P}\!\mathit{ic}}
\def\Picardfunctor{\mathrm{Pic}}
\def\Deformationcategory{\mathcal{D}\!\mathit{ef}}
\usepackage{cleveref}
\crefname{originaltheorem}{original theorem}{original theorems}
\Crefname{originaltheorem}{Original theorem}{Original theorems}
\makeatletter
\newlabel{theorem-vanishing-affine-curves}{{03SC}{source}{}{sourcepointer-1}{}}
\newlabel{theorem-vanishing-affine-curves@cref}{{[originaltheorem][0][]03SC}{[1][0][]source}{}{}{}}
\makeatother
\title{Cohomological dimension of a one-dimensional affine scheme over a field}
\author{The Stacks Project Authors}
\date{stacks-project:etale-cohomology.tex:local-snapshot-20261011; GNU FDL 1.2 or later}
\begin{document}
\maketitle
\noindent\textit{Editorial note.} All delivered mathematical text is the author's own, copied from the pinned Stacks Project chapter file etale-cohomology.tex; no statement or proof step was rewritten, reordered or removed.\par\medskip
\noindent\textit{Editorial note.} The proof is the author's complete argument; the final numerical inference from the vanishing region of the spectral sequence to the bound cd(X) at most 1 + cd(K) is left implicit by the author and is spelled out in this card's core\_obligations without altering the counted text.\par\medskip
\noindent\textit{Editorial note.} The label lemma-cd-curve-over-field occurs in no other staged Stacks chapter, so no cross-chapter restatement had to be weighed for this card.\par\medskip
\noindent\textit{Editorial note.} The single labelled reference occurring inside the counted spans is carried in assembly.external\_references as a non-load-bearing navigation pointer to the same chapter.\par\medskip
\noindent\textit{Editorial note.} The preamble is the chapter's own native preamble excerpt already used for this source: the Xy-pic setup (native preamble lines 14-18) and the native theorem declarations with the mathematical macros (native preamble lines 154-217).\par\medskip
\noindent\textit{Allowed standard prerequisites (editorial).}\begin{itemize}
\item Cohomological dimension cd(X) of a scheme. Same chapter, the definition of the cohomological dimension as the supremum of the degrees in which torsion sheaves can have nonzero cohomology. Supplies the quantity bounded in the statement.
\item Leray spectral sequence of a morphism of schemes. Same chapter, the Leray spectral sequence E\_2\textasciicircum{}\{p,q\} = H\textasciicircum{}p(S, R\textasciicircum{}q f\_*F) converging to H\textasciicircum{}\{p+q\}(X, F). Supplies the spectral sequence with which the proof reduces the bound to the vanishing of the higher direct images.
\item Vanishing of cohomology of torsion sheaves on an affine curve in degrees at least two. Same chapter, the theorem that H\textasciicircum{}i of a torsion sheaf on a one-dimensional affine scheme over a field vanishes for i at least 2. Supplies the vanishing of the stalks of R\textasciicircum{}q f\_*F in degrees at least 2, which is the input to the spectral sequence argument.
\end{itemize}\medskip
\ifdefined\hypertarget\hypertarget{sourcepointer-1}{}\fi
\noindent\textit{Source reference.} Original theorem 03SC: sources/stacks-etale-cohomology/pinned-original-source.tex, source lines 15656--15656. Same-chapter navigation pointer invoked at proof line 19668-19669, where the author uses that the cohomology of an affine curve vanishes in degrees at least two. The pointer only resolves that reference for the reader; the cited theorem is not reproduced inside the counted span, and no bibliography citation is involved.\par\medskip
\begin{lemma}
\label{lemma-cd-curve-over-field}
Let $K$ be a field. Let $X$ be a $1$-dimensional
affine scheme of finite type over $K$. Then
$\text{cd}(X) \leq 1 + \text{cd}(K)$.
\end{lemma}

\begin{proof}
Let $\mathcal{F}$ be an abelian torsion sheaf on $X_\etale$.
Consider the Leray spectral sequence for the morphism
$f : X \to \Spec(K)$. We obtain
$$
E_2^{p, q} = H^p(\Spec(K), R^qf_*\mathcal{F})
$$
converging to $H^{p + q}_\etale(X, \mathcal{F})$.
The stalk of $R^qf_*\mathcal{F}$ at a geometric point
$\Spec(\overline{K}) \to \Spec(K)$ is the cohomology of the
pullback of $\mathcal{F}$ to $X_{\overline{K}}$.
Hence it vanishes in degrees $\geq 2$ by
Theorem \ref{theorem-vanishing-affine-curves}.
\end{proof}

\end{document}
